%% ============================================================
%%  SECCIÓN 3 — Objetivos
%%  Responsable: Vera Vivanco, Jesús Rodolfo
%% ============================================================

\section{Objetivos}

\subsection{Objetivo general}

Determinar experimentalmente los coeficientes de dilatación lineal ($\alpha$) de muestras de tubos de aluminio, cobre y vidrio borosilicato, mediante el calentamiento controlado por inyección de vapor de agua.

\subsection{Objetivos específicos}

\begin{itemize}
    \item Medir la longitud inicial ($L_0$) de cada tubo a temperatura ambiente utilizando instrumentos de precisión.
    \item Cuantificar el alargamiento térmico lineal ($\Delta L$) de los materiales mediante el método de la aguja giratoria por rodadura sin deslizamiento.
    \item Registrar el ángulo de giro ($\Delta\theta$) que experimenta la aguja indicadora para calcular indirectamente el desplazamiento lineal del tubo.
    \item Propagar las incertidumbres instrumentales asociadas a las magnitudes físicas medidas para estimar la incertidumbre del coeficiente de dilatación.
    \item Comparar los coeficientes empíricos obtenidos con los valores de referencia teóricos y discutir las posibles fuentes de error sistemático y aleatorio.
\end{itemize}